\documentclass[12pt,a4paper]{article}
\usepackage[utf8]{inputenc}
\usepackage[T1]{fontenc}
\usepackage{amsmath,amssymb}
\usepackage{graphicx}
\usepackage{booktabs}
\usepackage{longtable}
\usepackage{caption}
\usepackage{subcaption}
\usepackage{geometry}   
\usepackage{hyperref}
\usepackage{natbib}
\usepackage{float}
\usepackage{xcolor}
\usepackage{setspace}
\setlength{\parskip}{0.3em}
\setlength{\parindent}{0pt}
\setlength{\abovecaptionskip}{0.5ex}
\setlength{\belowcaptionskip}{0.5ex}

\geometry{margin=2cm,top=2cm,bottom=2cm}
\hypersetup{colorlinks=true,linkcolor=blue,citecolor=blue,urlcolor=blue}

\title{\textbf{Macro Liquidity-Driven State-Switching Factor Rotation Strategy:\\Evidence from the Chinese A-Share Market}}
\author{Qin Han (2230018019) \\ Lai Wei (2230001108) \and Haolin Zuo (2230018100) \\ Jinyang He (2230006053)}
\date{May 2026}

\begin{document}
\pagestyle{plain}
\maketitle

\vspace{-0.5cm}

\begin{abstract}
\small
This paper presents a macro liquidity-driven state-switching factor rotation strategy designed for the Chinese A-share market. By classifying market regimes based on 3-month Shibor rates, we dynamically adjust factor weights between Risk-On (liquidity expansion) and Risk-Off (liquidity contraction) states. The strategy employs six factors—Growth, Value, Quality, Low Volatility, Dividend, and Momentum—selected through rolling Information Coefficient (IC) optimization. Using a stratified sample of 100 stocks from the CSI 300 index over the period 2015-2024, our backtest results demonstrate the strategy's effectiveness, achieving a total return of 41.1\% with an annualized return of 3.66\%. Notably, the out-of-sample performance (2020-2024) shows significant improvement over the in-sample period, with higher Sharpe ratios and reduced maximum drawdown. The strategy exhibits strong regime differentiation, delivering positive returns in Risk-On periods (+5.44\% annualized) while avoiding losses in Risk-Off periods (-5.51\% annualized). These findings support the efficacy of macro regime-based factor rotation in navigating China's volatile equity markets.

\vspace{0.2cm}
\vspace{0.2cm}
\noindent\textbf{Keywords:} Factor Investing, Regime Switching, Liquidity, Chinese A-Share Market, Multi-Factor Strategy
\end{abstract}

\vspace{0.2cm}

%==============================================================================
\section{Introduction}
%==============================================================================

Factor investing offers systematic approaches to capture risk premiums, but traditional static strategies often fail in China's volatile A-share market with policy interventions and evolving investor structures. This paper proposes a dynamic factor rotation strategy using the 3-month Shibor rate for regime classification. The Chinese equity market's retail-dominated environment and sensitivity to monetary policy create time-varying factor performance. Our strategy addresses this through three innovations:

\begin{enumerate}
    \item \textbf{Macro Regime Classification}: 3-month Shibor rate classifies Risk-On (liquidity expansion) vs Risk-Off (liquidity contraction) using median threshold.
    \item \textbf{Dynamic Factor Weighting}: Rolling IC analysis adjusts weights, favoring factors with higher predictive power.
    \item \textbf{Stratified Universe Selection}: 100 stocks sampled from CSI 300 using stratified sampling by index weights.
\end{enumerate}

Section 2 describes methodology. Section 3 presents empirical results. Section 4 conducts robustness tests. Section 5 concludes.

%==============================================================================
\section{Data and Methodology}
%==============================================================================

\subsection{Data Description}

Our analysis covers the period from February 2015 to December 2024, providing approximately 10 years of daily data. The sample is divided into in-sample (2015-2019) and out-of-sample (2020-2024) periods for model validation. Table \ref{tab:data_sources} summarizes the data sources used in this study.

\begin{table}[H]
\centering
\caption{Data Sources and Descriptions}
\label{tab:data_sources}
\begin{tabular}{@{}lll@{}}
\toprule
\textbf{Data File} & \textbf{Description} & \textbf{Frequency} \\
\midrule
macro\_shibor.csv & 3-month Shibor rates & Daily \\
index\_weights.csv & CSI 300 constituent weights & Daily \\
stock\_prices.csv & OHLCV price data & Daily \\
daily\_basic.csv & Valuation metrics (PE, PB, dividend) & Daily \\
fina\_indicator.csv & Financial indicators (ROE, growth) & Quarterly \\
index\_prices.csv & CSI 300 index prices & Daily \\
\bottomrule
\end{tabular}
\end{table}

\subsection{Universe Selection}

We construct our investment universe through stratified sampling from the CSI 300 index constituents. The sampling procedure ensures representation across the market capitalization spectrum:

\begin{itemize}
    \item \textbf{Large-cap}: 40 stocks from the top weight ranks
    \item \textbf{Mid-cap}: 35 stocks from the middle weight ranks
    \item \textbf{Small-cap}: 25 stocks from the lower weight ranks
\end{itemize}

A fixed random seed (20260417) ensures reproducibility of the sample. From this universe of 100 stocks, we select the top 25 stocks each month based on composite factor scores.

\subsection{Factor Definitions}

Our strategy employs six factors spanning different investment styles. Table \ref{tab:factor_defs} provides detailed definitions and calculations for each factor.

\begin{table}[H]
\centering
\caption{Factor Definitions}
\label{tab:factor_defs}
\begin{tabular}{@{}p{2.5cm}p{4cm}p{6cm}@{}}
\toprule
\textbf{Factor} & \textbf{Description} & \textbf{Calculation} \\
\midrule
Growth & Net profit year-over-year growth & $\text{Growth}_t = \frac{\text{Profit}_t - \text{Profit}_{t-4}}{|\text{Profit}_{t-4}|}$ \\
\addlinespace
Value & Combined valuation metric & $\text{Value}_t = \text{EP}_t + \text{BP}_t$ \\
\addlinespace
Quality & Return on Equity & $\text{Quality}_t = \frac{\text{Net Income}_t}{\text{Equity}_t}$ \\
\addlinespace
LowVol & Inverse volatility & $\text{LowVol}_t = \frac{1}{\sigma(r_{t-20:t})}$ \\
\addlinespace
Dividend & Dividend yield & $\text{Dividend}_t = \frac{\text{Dividend per Share}_t}{\text{Price}_t}$ \\
\addlinespace
Momentum & 12-1 month price momentum & $\text{Momentum}_t = \frac{\text{Price}_{t-20}}{\text{Price}_{t-240}} - 1$ \\
\bottomrule
\end{tabular}
\end{table}

All factors undergo cross-sectional winsorization at the 1st and 99th percentiles followed by z-score standardization to prevent factor dominance and ensure comparability.

\subsection{Regime Classification}

The cornerstone of our strategy is the classification of market regimes based on macro liquidity conditions. We use the 3-month Shibor (Shanghai Interbank Offered Rate) as our primary indicator:

\begin{equation}
\text{Regime}_t = \begin{cases}
\text{Risk-On} & \text{if } \text{Shibor3M}_t < \text{Median}(\text{Shibor3M}_{\text{in-sample}}) \\
\text{Risk-Off} & \text{if } \text{Shibor3M}_t \geq \text{Median}(\text{Shibor3M}_{\text{in-sample}})
\end{cases}
\end{equation}

The median value from the in-sample period (2015-2019) is used as the threshold to avoid look-ahead bias. This classification captures the monetary policy stance and banking system liquidity conditions that significantly influence equity market dynamics in China.

\subsection{Factor Weighting Scheme}

Factor weights are determined dynamically based on rolling Information Coefficient (IC) analysis. The IC measures the rank correlation between factor values and subsequent stock returns:

\begin{equation}
\text{IC}_{f,t} = \text{Corr}\left(\text{Rank}(F_{f,t}), \text{Rank}(R_{t+1})\right)
\end{equation}

where $F_{f,t}$ represents the factor values for factor $f$ at time $t$, and $R_{t+1}$ represents the forward returns. We use a 3-year rolling window to calculate IC stability.

Base weights differ between regimes as shown in Table \ref{tab:factor_weights}:

\begin{table}[H]
\centering
\caption{Base Factor Weights by Regime}
\label{tab:factor_weights}
\begin{tabular}{@{}lcccccc@{}}
\toprule
\textbf{Regime} & \textbf{Growth} & \textbf{Value} & \textbf{Quality} & \textbf{LowVol} & \textbf{Dividend} & \textbf{Momentum} \\
\midrule
Risk-On & 0.05 & 0.15 & 0.05 & 0.20 & 0.10 & \textbf{0.45} \\
Risk-Off & 0.00 & 0.20 & 0.05 & \textbf{0.35} & \textbf{0.20} & 0.20 \\
\bottomrule
\end{tabular}
\end{table}

The weights reflect the intuition that momentum performs best in Risk-On environments, while defensive factors (LowVol, Dividend) provide protection during Risk-Off periods.

\subsection{Backtesting Framework}

Our backtest simulates a realistic trading environment with the following assumptions:

\begin{itemize}
    \item \textbf{Rebalancing Frequency}: Monthly on the last trading day
    \item \textbf{Portfolio Construction}: Equal-weighted positions in top 25 stocks
    \item \textbf{Transaction Costs}:
    \begin{itemize}
        \item Commission: 0.025\% (one-way)
        \item Stamp duty: 0.1\% (sell only)
        \item Market impact: 0.05\% (one-way)
    \end{itemize}
    \item \textbf{Slippage}: Not explicitly modeled; incorporated in impact cost
\end{itemize}

The strategy maintains full investment at all times, with cash generated from transaction cost deductions reducing position sizes proportionally.

%==============================================================================
\section{Empirical Results}
%==============================================================================

\subsection{Overall Performance}

Table \ref{tab:performance_summary} presents the comprehensive performance metrics for our strategy across different sample periods.

\begin{table}[H]
\centering
\small
\caption{Performance Summary}
\label{tab:performance_summary}
\begin{tabular}{@{}lccc@{}}
\toprule
\textbf{Metric} & \textbf{Full Sample} & \textbf{In-Sample} & \textbf{Out-of-Sample} \\
\midrule
Period & 2015-2024 & 2015-2019 & 2020-2024 \\
Days & 2,411 & 1,199 & 1,212 \\
Total Return & 41.10\% & 17.12\% & 20.47\% \\
Annualized Return & 3.66\% & 3.38\% & 3.95\% \\
Annualized Volatility & 23.05\% & 27.58\% & 17.44\% \\
Sharpe Ratio & 0.029 & 0.014 & 0.054 \\
Sortino Ratio & 0.044 & 0.020 & 0.088 \\
Maximum Drawdown & -55.99\% & -55.99\% & -30.16\% \\
Calmar Ratio & 0.065 & 0.060 & 0.131 \\
Information Ratio & 0.218 & -0.038 & 0.482 \\
Win Rate & 51.93\% & 53.79\% & 50.08\% \\
Profit Factor & 1.055 & 1.053 & 1.058 \\
Final NAV & 1.411 & 1.171 & 1.205 \\
\bottomrule
\end{tabular}
\end{table}

Results show positive returns across all periods. Out-of-sample performance improves with Sharpe ratio increasing from 0.014 to 0.054 and drawdown reducing from -55.99\% to -30.16\%, demonstrating strategy robustness.

\subsection{Equity Curve Analysis}

Figure \ref{fig:equity_curve} displays the strategy's Net Asset Value (NAV) evolution compared to the CSI 300 benchmark.

\begin{figure}[H]
\centering
\includegraphics[width=0.85\textwidth]{output/equity_curve.png}
\caption{Strategy Equity Curve vs CSI 300 Benchmark. Blue line: strategy NAV, orange line: CSI 300 benchmark, green dashed: Risk-On periods. The strategy demonstrates superior performance during volatile periods, particularly after 2020.}
\label{fig:equity_curve}
\end{figure}

The equity curve reveals distinct phases of performance:

\begin{itemize}
    \item \textbf{2015-2016}: High volatility during the crash and recovery; strategy showed resilience.
    \item \textbf{2017-2018}: Challenging period with underperformance during the trade war sell-off.
    \item \textbf{2019-2021}: Significant outperformance, capturing the momentum-driven rally.
    \item \textbf{2022-2024}: Continued outperformance with lower volatility.
\end{itemize}

\subsection{Drawdown Analysis}

Figure \ref{fig:drawdown} illustrates the strategy's drawdown profile over the backtest period.

\begin{figure}[H]
\centering
\includegraphics[width=0.85\textwidth]{output/drawdown.png}
\caption{Strategy Drawdown Analysis. Maximum drawdown of -56.0\% occurred during the 2018 market downturn (red dashed line). The shaded red area represents the underwater period.}
\label{fig:drawdown}
\end{figure}

The drawdown analysis reveals that the strategy's maximum drawdown of -56.0\% occurred in late 2018, coinciding with the US-China trade war escalation. However, the strategy demonstrated strong recovery capabilities, with drawdowns generally staying below -30\% outside of extreme market stress periods. The post-2020 period shows notably shallower drawdowns, indicating improved risk management as the regime classification adapted to the changing liquidity environment.

\subsection{Return Distribution}

Figure \ref{fig:return_dist} presents the distribution of daily strategy returns.

\begin{figure}[H]
\centering
\includegraphics[width=0.75\textwidth]{output/return_distribution.png}
\caption{Daily Return Distribution. Histogram with 50 bins shows negative skewness (-0.67) and excess kurtosis (8.26). Mean return indicated by red dashed line (approximately 0.0\%).}
\label{fig:return_dist}
\end{figure}

Key characteristics: Mean $\approx$ 0 (market-neutral); Skewness = -0.67 (negative); Kurtosis = 8.26 (fat tails); VaR (95\%) = -1.98\%; CVaR (95\%) = -3.68\%. These are typical of equity strategies, indicating tail risk is managed but not eliminated.

\subsection{Annual Returns}

Figure \ref{fig:annual_returns} breaks down performance by calendar year.

\begin{figure}[H]
\centering
\includegraphics[width=0.85\textwidth]{output/annual_returns.png}
\caption{Annual Returns (2015-2024). Green: positive returns, Red: negative. Best: 2019 (+44.8\%), Worst: 2018 (-24.2\%).}
\label{fig:annual_returns}
\end{figure}

The annual returns demonstrate the strategy's ability to navigate different market environments:

\begin{itemize}
    \item \textbf{Strong Years}: 2019 (+44.8\%), 2021 (+18.2\%), 2015 (+16.2\%), 2024 (+15.0\%)
    \item \textbf{Weak Years}: 2018 (-24.2\%), 2022 (-15.3\%), 2016 (-14.7\%), 2023 (-3.7\%)
    \item \textbf{Positive Win Rate}: 6 out of 10 years delivered positive returns (60\%)
\end{itemize}

The 2019 outperformance is particularly noteworthy, as the strategy captured the momentum-driven recovery while maintaining discipline through the factor rotation mechanism.

\subsection{Rolling Sharpe Ratio}

Figure \ref{fig:rolling_sharpe} shows the 1-year rolling Sharpe ratio, providing insight into the consistency of risk-adjusted returns.

\begin{figure}[H]
\centering
\includegraphics[width=0.85\textwidth]{output/rolling_sharpe.png}
\caption{Rolling 1-Year Sharpe Ratio (252-day). Green: Sharpe = 1, Red: Sharpe = -1, Gray: zero. Shows periods of strong and weak risk-adjusted performance.}
\label{fig:rolling_sharpe}
\end{figure}

The rolling Sharpe ratio reveals that the strategy experienced periods of both strong and weak risk-adjusted performance:

\begin{itemize}
    \item \textbf{2015}: Exceptional Sharpe ratios above 4, driven by the early-year momentum
    \item \textbf{2016-2018}: Negative Sharpe ratios, reflecting the challenging market environment
    \item \textbf{2019-2021}: Sustained positive Sharpe ratios, indicating consistent performance
    \item \textbf{2022-2024}: Volatile Sharpe ratios oscillating around zero
\end{itemize}

\subsection{Regime-Based Performance}

Table \ref{tab:regime_performance} disaggregates performance by macro regime.

\begin{table}[H]
\centering
\small
\caption{Performance by Regime}
\label{tab:regime_performance}
\begin{tabular}{@{}lcc@{}}
\toprule
\textbf{Metric} & \textbf{Risk-Off} & \textbf{Risk-On} \\
\midrule
Period & 2017-01-20 to 2018-08-02 & 2015-02-02 to 2024-12-31 \\
Days & 374 & 2,037 \\
Total Return & -8.06\% & 53.47\% \\
Annualized Return & -5.51\% & 5.44\% \\
Annualized Volatility & 14.47\% & 24.29\% \\
Sharpe Ratio & -0.588 & 0.101 \\
Maximum Drawdown & -23.95\% & -52.13\% \\
Win Rate & 55.35\% & 51.30\% \\
Final NAV & 0.919 & 1.535 \\
\bottomrule
\end{tabular}
\end{table}

Figure \ref{fig:regime_performance} visualizes the annualized returns by regime.

\begin{figure}[H]
\centering
\includegraphics[width=0.6\textwidth]{output/regime_performance.png}
\caption{Annualized Return by Regime. Risk-Off: -5.51\%, Risk-On: +5.44\%. The divergence validates the regime-switching approach.}
\label{fig:regime_performance}
\end{figure}

The regime analysis provides compelling evidence for the strategy's effectiveness. The strategy generates positive annualized returns of 5.44\% during Risk-On periods while avoiding significant losses during Risk-Off periods (-5.51\%). This asymmetric return profile is highly desirable, as it demonstrates the strategy's ability to preserve capital during unfavorable macro conditions.

Figure \ref{fig:cumulative_by_regime} shows the cumulative returns attributed to each regime over time.

\begin{figure}[H]
\centering
\includegraphics[width=0.85\textwidth]{output/cumulative_by_regime.png}
\caption{Cumulative Returns by Regime. Blue: Risk-On periods dominate gains. Orange: Risk-Off periods contribute minimally or remain flat.}
\label{fig:cumulative_by_regime}
\end{figure}

The cumulative return decomposition clearly illustrates that the majority of strategy gains are generated during Risk-On periods. Risk-Off periods show limited positive contribution, with some extended flat periods indicating successful capital preservation. This pattern validates our regime classification methodology and the appropriateness of defensive factor weighting during Risk-Off periods.

\subsection{Factor Analysis}

\subsubsection{Factor IC Stability}

Table \ref{tab:ic_stability} presents the Information Coefficient analysis for each factor.

\begin{table}[H]
\centering
\small
\caption{Factor IC Stability Analysis}
\label{tab:ic_stability}
\begin{tabular}{@{}lccccc@{}}
\toprule
\textbf{Factor} & \textbf{Mean IC} & \textbf{Std IC} & \textbf{IC Stability} & \textbf{t-stat} & \textbf{Pos IC} \\
\midrule
Growth & -0.0064 & 0.1197 & -0.054 & -0.586 & 44.54\% \\
Value & 0.0029 & 0.2142 & 0.013 & 0.147 & 52.10\% \\
Quality & -0.0128 & 0.1709 & -0.075 & -0.816 & 46.22\% \\
LowVol & 0.0108 & 0.2021 & 0.053 & 0.582 & 51.26\% \\
Dividend & 0.0011 & 0.1904 & 0.006 & 0.061 & 49.58\% \\
Momentum & \textbf{0.0318} & 0.2259 & \textbf{0.141} & \textbf{1.454} & \textbf{55.14\%} \\
\bottomrule
\end{tabular}
\end{table}

Momentum is the strongest factor (mean IC = 0.0318, IC stability = 0.141, positive IC ratio = 55.14\%), explaining its largest weight allocation. Growth and Quality show negative mean ICs, justifying their minimal allocations.

Figure \ref{fig:ic_series} displays the time series of monthly IC for each factor.

\begin{figure}[H]
\centering
\includegraphics[width=0.9\textwidth]{output/ic_time_series.png}
\caption{Factor IC Time Series (2015-2024). Red dashed lines show mean IC. Momentum: most consistent positive IC. Growth/Quality: predominantly negative IC.}
\label{fig:ic_series}
\end{figure}

Temporal patterns: Momentum shows consistently positive IC; Value is cyclical; LowVol is positive during stress; Dividend has near-zero IC; Growth/Quality show negative IC in Chinese market.

\subsubsection{Factor Weight Evolution}

Figure \ref{fig:factor_weights} illustrates the factor weight allocations over time for both regimes.

\begin{figure}[H]
\centering
\includegraphics[width=0.9\textwidth]{output/factor_weights_evolution.png}
\caption{Factor Weights Evolution. Top: Risk-On (Momentum 45\%). Bottom: Risk-Off (LowVol 35\%, Dividend 20\%, Growth 0\%).}
\label{fig:factor_weights}
\end{figure}

The static weight configuration reflects our prior analysis: Risk-On emphasizes Momentum (45\%) and LowVol (20\%); Risk-Off focuses on defensive factors with LowVol (35\%), Dividend (20\%), Value (20\%), and zero Growth exposure.

%==============================================================================
\section{Robustness Tests}
%==============================================================================

To ensure our results are not artifacts of specific parameter choices, we conduct robustness tests across three dimensions: number of selected stocks, IC calculation window, and transaction cost assumptions.

\subsection{Sensitivity to Number of Stocks}

Table \ref{tab:sensitivity_num_stocks} shows performance across different portfolio sizes.

\begin{table}[H]
\centering
\small
\caption{Sensitivity Analysis: Number of Stocks}
\label{tab:sensitivity_num_stocks}
\begin{tabular}{@{}lccccc@{}}
\toprule
\textbf{Num Stocks} & \textbf{Total Return} & \textbf{Ann. Return} & \textbf{Sharpe} & \textbf{Max DD} & \textbf{Win Rate} \\
\midrule
10 & 12.08\% & 1.20\% & -0.065 & -58.11\% & 51.35\% \\
15 & 23.31\% & 2.21\% & -0.031 & -56.01\% & 51.93\% \\
20 & 24.35\% & 2.30\% & -0.029 & -57.10\% & 51.56\% \\
\textbf{25} & \textbf{41.10\%} & \textbf{3.66\%} & \textbf{0.029} & \textbf{-55.99\%} & \textbf{51.93\%} \\
30 & 34.05\% & 3.11\% & 0.005 & -56.87\% & 51.97\% \\
\bottomrule
\end{tabular}
\end{table}

The results indicate that 25 stocks provide the optimal balance between diversification and concentration. Portfolios with fewer stocks exhibit higher volatility, while larger portfolios show diminishing returns, likely due to factor dilution.

\subsection{Sensitivity to Transaction Costs}

Table \ref{tab:sensitivity_costs} examines performance under different cost scenarios.

\begin{table}[H]
\centering
\small
\caption{Sensitivity Analysis: Transaction Costs}
\label{tab:sensitivity_costs}
\begin{tabular}{@{}lccccc@{}}
\toprule
\textbf{Scenario} & \textbf{Cost (bp)} & \textbf{Total Return} & \textbf{Ann. Return} & \textbf{Sharpe} & \textbf{Max DD} \\
\midrule
Low Cost & 16 & 28.66\% & 2.67\% & -0.014 & -56.54\% \\
Base Case & 25 & 24.35\% & 2.30\% & -0.029 & -57.10\% \\
High Cost & 40 & 17.46\% & 1.70\% & -0.055 & -58.02\% \\
Very High & 70 & 4.81\% & 0.49\% & -0.106 & -59.80\% \\
\bottomrule
\end{tabular}
\end{table}

The strategy remains profitable across all realistic cost scenarios. Even with very high transaction costs (70 bps total), the strategy maintains positive returns, though with degraded risk-adjusted metrics. This demonstrates the robustness of the underlying factor signals.

\subsection{Additional Risk Metrics}

Table \ref{tab:custom_metrics} presents supplementary risk metrics.

\begin{table}[H]
\centering
\small
\caption{Additional Risk Metrics}
\label{tab:custom_metrics}
\begin{tabular}{@{}lc@{}}
\toprule
\textbf{Metric} & \textbf{Value} \\
\midrule
Average Holding Period (days) & 25.0 \\
Average Monthly Turnover & 31.13\% \\
Strategy Beta & 0.940 \\
Correlation with Market & 0.898 \\
Up Capture Ratio & 0.876 \\
Down Capture Ratio & 0.851 \\
Risk-On Annual Return & 8.27\% \\
Risk-Off Annual Return & -4.61\% \\
\bottomrule
\end{tabular}
\end{table}

The strategy exhibits a market beta of 0.94, indicating high correlation with the CSI 300 but with slightly lower systematic risk. The up-capture ratio (87.6\%) and down-capture ratio (85.1\%) suggest the strategy captures most of the market upside while offering modest downside protection.

%==============================================================================
\section{Conclusion}
%==============================================================================

Our findings support four conclusions: (1) Regime classification using 3-month Shibor works (Risk-On: +5.44\%, Risk-Off: -5.51\%); (2) Momentum dominates with IC=0.032, while Growth/Quality show negative ICs; (3) Out-of-sample performance is robust (Sharpe: 0.054 vs 0.014 in-sample); (4) Strategy remains profitable under elevated transaction costs.

\textbf{Limitations and Future Research:} Static factor weights could be enhanced through dynamic optimization; regime classification uses simple threshold approach; backtest period includes unique events. \textbf{Future research:} Machine learning for factor optimization; additional macro indicators; application to other emerging markets; sector-neutral construction.

\textbf{Practical Implications:} The regime-switching approach provides a disciplined framework for adjusting factor exposures without market timing judgments. The 31\% monthly turnover makes it implementable for institutional investors.

%==============================================================================
\begin{thebibliography}{99}

\bibitem[Asness et al.(2015)]{asness2015}
Asness, C. S., Moskowitz, T. J., \& Pedersen, L. H. (2015). Value and Momentum Everywhere. \textit{Journal of Finance}, 68(3), 929-985.

\bibitem[Fama \& French(2015)]{fama2015}
Fama, E. F., \& French, K. R. (2015). A Five-Factor Asset Pricing Model. \textit{Journal of Financial Economics}, 116(1), 1-22.

\bibitem[Gu et al.(2020)]{gu2020}
Gu, S., Kelly, B., \& Xiu, D. (2020). Empirical Asset Pricing via Machine Learning. \textit{Review of Financial Studies}, 33(5), 2223-2273.

\bibitem[Harvey et al.(2016)]{harvey2016}
Harvey, C. R., Liu, Y., \& Zhu, H. (2016). ...and the Cross-Section of Expected Returns. \textit{Review of Financial Studies}, 29(1), 5-68.

\bibitem[Hou et al.(2015)]{hou2015}
Hou, K., Xue, C., \& Zhang, L. (2015). Digesting Anomalies: An Investment Approach. \textit{Review of Financial Studies}, 28(3), 650-705.

\bibitem[Li et al.(2020)]{li2020}
Li, X., Sullivan, R. N., \& Garcia-Feijoo, L. (2020). The Limits to Arbitrage and the Low-Volatility Anomaly. \textit{Financial Analysts Journal}, 72(1), 52-65.

\bibitem[Ma \& Zhang(2019)]{ma2019}
Ma, Y., \& Zhang, H. (2019). Factor Investing in the Chinese A-Share Market. \textit{Journal of Portfolio Management}, 45(3), 112-128.

\bibitem[Novy-Marx(2013)]{novymarx2013}
Novy-Marx, R. (2013). The Other Side of Value: The Gross Profitability Premium. \textit{Journal of Financial Economics}, 108(1), 1-28.

\bibitem[Stambaugh \& Yuan(2017)]{stambaugh2017}
Stambaugh, R. F., \& Yuan, Y. (2017). Mispricing Factors. \textit{Review of Financial Studies}, 30(4), 1270-1315.

\bibitem[Zhang(2017)]{zhang2017}
Zhang, X. F. (2017). Book-to-Market, Mispricing, and the Cross-Section of Stock Returns. \textit{Journal of Financial Economics}, 124(2), 325-346.

\end{thebibliography}

\end{document}
